\documentclass[parskip=full]{scrartcl}
\usepackage[top=2.5cm, bottom=2.5cm, left=2.5cm, right=2.5cm]{geometry}
\usepackage[utf8]{inputenc}
\usepackage[T1]{fontenc}
\usepackage[german]{babel}
\usepackage{hyperref}
\usepackage[toc, nonumberlist]{glossaries}
\usepackage{graphicx}
\usepackage{enumitem}
\usepackage{float}
\usepackage{color}

\hypersetup{
  colorlinks=false,
  linktoc=all,
  hidelinks,
}

\title{OSIP - Implementierung}
\subtitle{OPC UA Simulator for Industrial Plants}
\author{
    M. Armbruster\\
    D. Kahles\\
    H. Lehmann\\
    M. Schwarzmann\\
    N. Wilhelm
}

\begin{document}
\maketitle
\thispagestyle{empty}
\vspace{20px}
\begin{center}
  \includegraphics[scale=0.4]{../icon.png}
\end{center}
\pagebreak
\tableofcontents
\pagebreak

\section{Einleitung}


\section{Änderungen zum Entwurf}


\section{Implementierte Muss- und Kannkriterien}


\section{Verzögerungen im Implementierungsplan}


\section{Probleme bei der Implementierung}
Im Pflichtenheft war geplant, die von OPC UA bereitgestellten Alarme zu benutzen, um bei einem ungewöhnlichen Zustand (z.B. Überlauf)
die Überwachungskonsole zu benachrichtigen. Leider unterstützt der Milo Server keine Alarme, siehe auch die entsprechende Antwort
des Milo-Entwicklers auf Stack Overflow: \href{http://stackoverflow.com/a/42161092/}{http://stackoverflow.com/a/42161092/}.
Somit mussten wir auch bei den Alarmen auf reguläre Subscriptions zurückgreifen, welche wir nun mit einem kürzeren Intervall
als die restlichen Produktionsvariablen abfragen.

\section{Unit tests}


\pagebreak
\phantomsection
\addcontentsline{toc}{section}{\listfigurename}
\listoffigures

\end{document}
